\documentclass[11pt]{article}
\usepackage[margin=1in]{geometry}
\usepackage{amsmath,amssymb,amsthm}
\theoremstyle{definition}
\newtheorem{theorem}{Theorem}[section]
\newtheorem{proposition}[theorem]{Proposition}
\title{Analysis 10: Fixed Crossing Delay and Concentration\\
for Weighted Rolling Validation}
\author{Aileen Tao}
\date{}
\begin{document}
\maketitle

\section{Motivation: from exact alignment to fixed crossing delay}

Analysis 09 showed that exact crossing alignment concentrates positive
rolling-validation (RV) weights on an increasingly small fraction of
recent validation times, obstructing Chebyshev concentration on the
moving-crossing scale. Analysis 10 studies the tradeoff between fixed
crossing delay and concentration:
\[
\begin{gathered}
\text{If we allow }\frac{n_{{\rm RV},w}}{n_{\rm pop}}\to\lambda>1,\\
\text{how much do averaging and concentration improve,}\\
\text{and is that enough for the Chebyshev ratio to vanish?}
\end{gathered}
\]
Analysis 06 established power-weight concentration against an oracle.
We reuse its full variance calculation, distinguishing that benchmark
from Analysis 09's moving-crossing model.

Let $Y_i=\theta+\xi_i$, with independent $\xi_i\sim N(0,1)$ and
$0<|\theta|<1$. Model 0 predicts with the training sample mean;
Model 1 predicts zero. Both use squared loss. At validation time
$i=2,\ldots,n+1$, the training size is $i-1$. Throughout, the weights
$w_i>0$ form a deterministic sequence independent of the data and
$\theta$. The parameter $\lambda$ is fixed and finite unless a
subsequent parameter limit is explicitly stated.

\section{General constant crossing ratio}

Analysis 09 gives
\[
\Delta_{{\rm RV},w}(n)
=\sum_{i=2}^{n+1}w_i\left(\frac1{i-1}-\theta^2\right),
\qquad
r_n=\frac{\sum_{i=2}^{n+1}w_i/(i-1)}{\sum_{i=2}^{n+1}w_i}.
\]
Dividing by positive total weight preserves the crossing:
\[
\frac{\Delta_{{\rm RV},w}(n)}{\sum_{i=2}^{n+1}w_i}
=r_n-\theta^2,\qquad \Delta_R(n)=\frac1n-\theta^2.
\]
Thus $r_n$ plays the same role for weighted RV that $1/n$ plays for
population risk. Write $n_{{\rm RV},w}$ and $n_{\rm pop}$ for the first
strict negative indices of these differences. Since $r_n\ge1/n$, the
normalized expected weighted-RV difference is always at least as large
as the population-risk difference. Thus weighted RV cannot cross earlier
than population risk:
\[
n_{{\rm RV},w}\ge n_{\rm pop}
\]
whenever the weighted crossing exists. Therefore any finite limiting
crossing ratio must satisfy
\[
\lambda\ge1.
\]

\begin{samepage}
\begin{theorem}[Constant crossing delay]\label{thm:crossing}
For fixed $\lambda\ge1$,
\[
nr_n\to\lambda
\quad\Longleftrightarrow\quad
r_n\sim\frac{\lambda}{n}.
\]
Under this condition, $n_{{\rm RV},w}$ exists for every
$0<|\theta|<1$, and
\[
n_{\rm pop}=\lfloor\theta^{-2}\rfloor+1\sim\theta^{-2},
\qquad
\frac{n_{{\rm RV},w}}{n_{\rm pop}}\to\lambda
\qquad(|\theta|\to0).
\]
\end{theorem}
\noindent\textbf{Interpretation.} The value $\lambda=1$ is the earliest
possible asymptotic crossing; $\lambda>1$ imposes a fixed delay.
\textbf{Proof.} See Appendix~\ref{app:bracketing}, extending Analysis 09,
Appendix B.4.
\end{samepage}

\section{Power weights and the $\lambda$-tradeoff}

Analysis 06 studied the same power-weight family $w_i=i^\alpha$ under
an oracle comparison. Here we use this family to quantify the crossing-delay
and averaging tradeoff in the present model.

\begin{proposition}[Power weights and crossing delay]\label{prop:power}
For $w_i=i^\alpha$ with fixed $\alpha>0$,
\[
\sum_{i=2}^{n+1}w_i\sim\frac{n^{\alpha+1}}{\alpha+1},
\qquad
\sum_{i=2}^{n+1}\frac{w_i}{i-1}\sim\frac{n^\alpha}{\alpha},
\qquad
\sum_{i=2}^{n+1}w_i^2\sim\frac{n^{2\alpha+1}}{2\alpha+1}.
\]
Consequently,
\[
r_n\sim\frac{\alpha+1}{\alpha n}
\quad\Longrightarrow\quad nr_n\to\frac{\alpha+1}{\alpha}.
\]
Hence, by Theorem~\ref{thm:crossing},
\[
\lambda=\frac{\alpha+1}{\alpha}=1+\frac1\alpha,\qquad
\frac{n_{{\rm RV},w}}{n_{\rm pop}}\to\lambda.
\]
\end{proposition}
\noindent\textbf{Interpretation.} The relation
\[
\lambda=1+\frac1\alpha
\]
shows how the power exponent controls the crossing delay. Equivalently,
\[
\alpha=\frac1{\lambda-1},
\]
so every fixed delay $\lambda>1$ can be achieved by an appropriate
choice of $\alpha$. For power weights,
\[
n\frac{\sum_{i=2}^{n+1}w_i^2}
{\left(\sum_{i=2}^{n+1}w_i\right)^2}
\to\frac{(\alpha+1)^2}{2\alpha+1}.
\]
Substituting $\alpha=1/(\lambda-1)$ gives
\[
\frac{(\alpha+1)^2}{2\alpha+1}=\frac{\lambda^2}{\lambda^2-1}.
\]
This normalized square-weight factor measures how concentrated the
weights are. Values closer to $1$ correspond to more evenly spread
weights and stronger averaging, while larger values correspond to
more concentrated weights and weaker averaging.

The following parameter limits are taken after $n\to\infty$ at fixed
$\lambda>1$. As $\lambda\downarrow1$,
\[
\lambda\downarrow1\quad\Longrightarrow\quad\alpha\to\infty,
\qquad \lambda^2-1\to0^+,
\qquad
\frac{\lambda^2}{\lambda^2-1}\to\infty.
\]
As $\lambda$ approaches $1$, $\alpha$ becomes large and the power
weights put more mass on the latest validation times. The weights are
therefore more concentrated, so averaging becomes weaker. This matches
Analysis 09: exact alignment ($\lambda=1$) forces the normalized
square-weight factor to diverge.

Conversely,
\[
\lambda\to\infty\quad\Longrightarrow\quad\alpha\downarrow0,
\qquad
\frac{\lambda^2}{\lambda^2-1}
=\frac1{1-\lambda^{-2}}\longrightarrow1.
\]
As $\lambda\to\infty$, $\alpha\to0$, so the power weights become close
to equal weights. The weights are more evenly spread, the normalized
square-weight factor approaches $1$, and averaging becomes stronger.
The price is a much later crossing.
\textbf{Proof.} See Appendix~\ref{app:power-sums} and
Theorem~\ref{thm:crossing}.

\section{Concentration on the moving-crossing scale for fixed $\lambda>1$}

Analysis 09 gives the realized-score lower bound
\[
\Xi_{n,w}=\sum_{i=2}^{n+1}w_i
\{(Y_i-\bar Y_{i-1})^2-Y_i^2\},
\qquad
\mathbb E\Xi_{n,w}=\left(\sum_{i=2}^{n+1}w_i\right)(r_n-\theta^2),
\]
\begin{equation}\label{eq:lower}
\frac{\operatorname{Var}(\Xi_{n,w})}{\{\mathbb E\Xi_{n,w}\}^2}
\ge\frac{4\theta^2\sum_{i=2}^{n+1}w_i^2}
{\left(\sum_{i=2}^{n+1}w_i\right)^2(r_n-\theta^2)^2},
\end{equation}
when the mean is nonzero (Analysis 09, Appendices D.1--D.2).

Under
\[
nr_n\to\lambda\in[1,\infty),\qquad
n\theta^2\to c\in(0,\infty),\qquad |\theta|\to0,
\]
we have
\[
r_n\sim\frac{\lambda}{n},\qquad \theta^2\sim\frac cn,
\qquad r_n-\theta^2\sim\frac{\lambda-c}{n}\quad(c\ne\lambda).
\]
Substitution into \eqref{eq:lower} gives
\[
\frac{\operatorname{Var}(\Xi_{n,w})}{\{\mathbb E\Xi_{n,w}\}^2}
\ge(1+o(1))\frac{4c}{(\lambda-c)^2}\,
n\frac{\sum_{i=2}^{n+1}w_i^2}
{\left(\sum_{i=2}^{n+1}w_i\right)^2}.
\]
For positive weights, Cauchy--Schwarz gives
\[
\left(\sum_{i=2}^{n+1}w_i\right)^2
\le n\sum_{i=2}^{n+1}w_i^2,
\]
so
\[
n\frac{\sum_{i=2}^{n+1}w_i^2}
{\left(\sum_{i=2}^{n+1}w_i\right)^2}\ge1.
\]
This removes dependence on the specific weights and gives the general
lower bound below.

\begin{samepage}
\begin{proposition}[Positive-weight lower bound]\label{prop:lower}
Under the preceding limits and positive deterministic weights,
\[
\liminf\frac{\operatorname{Var}(\Xi_{n,w})}
{\{\mathbb E\Xi_{n,w}\}^2}
\ge\frac{4c}{(\lambda-c)^2}>0,\qquad c\ne\lambda.
\]
At $c=\lambda$, the leading normalized mean vanishes; the ratio diverges
when the mean is nonzero and is undefined at exact mean-zero points.
\end{proposition}
\noindent\textbf{Interpretation.} For any finite crossing delay $\lambda$,
the Chebyshev ratio does not vanish on this scale.
\textbf{Proof.} See Appendix~\ref{app:general-lower}.
\end{samepage}

\medskip
For power weights $w_i=i^\alpha$, $\alpha>0$, the current-model mean is
\[
\mathbb E\Xi_{n,w}
=\left(\sum_{i=2}^{n+1}w_i\right)(r_n-\theta^2)
\sim\frac{n^{\alpha+1}}{\alpha+1}\frac{\lambda-c}{n}.
\]
Since $\lambda=(\alpha+1)/\alpha$,
\[
\mathbb E\Xi_{n,w}
\sim n^\alpha\left(\frac1\alpha-\frac{c}{\alpha+1}\right),
\qquad c\ne\lambda.
\]
The quadratic-noise component is unchanged from Analysis 06, so its
variance asymptotic can be reused directly:
\[
\operatorname{Var}(\Xi_n^{(\alpha)})\sim\frac{2n^{2\alpha}}{\alpha+1}.
\]
The present model adds the linear-noise term
\[
-2\theta\sum_{i=2}^{n+1}i^\alpha\xi_i,
\]
whose variance is
\[
4\theta^2\sum_{i=2}^{n+1}i^{2\alpha}
\sim\frac{4c}{2\alpha+1}n^{2\alpha}.
\]
Therefore,
\[
\operatorname{Var}(\Xi_{n,w})
\sim n^{2\alpha}\left(\frac2{\alpha+1}+\frac{4c}{2\alpha+1}\right).
\]
The first term is the Analysis 06 quadratic component; the second is
the additional linear-noise contribution (Appendix~\ref{app:transfer}).

Combining the mean and variance gives the following result.
\begin{theorem}[Full moving-crossing ratio for power weights]\label{thm:ratio}
For fixed $\alpha>0$, $\lambda=1+1/\alpha$, and the preceding
moving-crossing limits with $c\ne\lambda$,
\begin{equation}\label{eq:full-ratio}
\frac{\operatorname{Var}(\Xi_{n,w})}{\{\mathbb E\Xi_{n,w}\}^2}
\to
\frac{2\lambda/(\lambda-1)+4c\lambda^2/(\lambda^2-1)}
{(\lambda-c)^2}\in(0,\infty).
\end{equation}
At $c=\lambda$, the ratio diverges when defined.
\end{theorem}
\noindent\textbf{Interpretation.} For fixed finite $\lambda>1$, the
limit is positive for $c\ne\lambda$, so the Chebyshev ratio does not
vanish; at $c=\lambda$, it diverges when defined.
\textbf{Proof.} See Appendices~\ref{app:transfer}--\ref{app:power-ratio}.

\section{Conclusion and implications}

Allowing a fixed crossing delay improves averaging compared with the
exact-alignment case studied in Analysis 09, because the weights do not
need to be as concentrated on the latest validation times.

However, any fixed finite delay is still not enough to make the
variance-to-mean-squared ratio vanish on the moving-crossing scale.
The ratio remains positive away from the weighted crossing and becomes
unbounded at the crossing itself, when defined.

The main tradeoff is clear: better crossing alignment requires more
concentrated weights and weaker averaging, while stronger averaging
requires accepting a larger crossing delay.

\appendix
\section*{Appendix}
\section{General crossing theorem}

\subsection{Population risk, positivity, and monotonicity}\label{app:monotonicity}
An independent new response has variance $1$. The sample-mean
predictor has estimation variance $1/n$ and no bias, while the zero
predictor has squared bias $\theta^2$. Thus their risks are $1+1/n$
and $1+\theta^2$, respectively. Their difference is negative exactly
when $n>\theta^{-2}$, giving
$n_{\rm pop}=\lfloor\theta^{-2}\rfloor+1$. This is the calculation
in Analysis 09, Appendix A.1, and applies to either sign of $\theta$.

The one-time expected validation difference is
$1/(i-1)-\theta^2$. Linearity of expectation gives the weighted mean
in Section 2 without requiring independence of validation losses.
Every reciprocal in $r_n$ is at least $1/n$, so $r_n\ge1/n$ and
the weighted strict crossing cannot precede the population crossing.
For $n>1$, the inequality is strict because the earlier terms have
positive weights.

Adding the next positive weight gives
\[
r_{n+1}=
\frac{(\sum_{i=2}^{n+1}w_i)r_n+w_{n+2}/(n+1)}
{\sum_{i=2}^{n+1}w_i+w_{n+2}}.
\]
Since $r_n\ge1/n>1/(n+1)$, this weighted average is strictly smaller
than $r_n$. Thus once $r_n<\theta^2$ occurs, it remains so.
Under $nr_n\to\lambda<\infty$, $r_n\to0$; hence a first strict
crossing exists for each fixed $\theta\ne0$. Also $r_1=1>\theta^2$.

\subsection{Epsilon bracketing of the first strict crossing}\label{app:bracketing}
Because $\lambda>0$, the identity
$r_n/(\lambda/n)=nr_n/\lambda$ proves the equivalence in
Theorem~\ref{thm:crossing}. To prove its crossing conclusion, put
$h=\theta^2\downarrow0$ and fix $\varepsilon\in(0,1)$. Set
\[
n_- =\left\lfloor\frac{(1-\varepsilon)\lambda}{h}\right\rfloor,
\qquad
n_+ =\left\lceil\frac{(1+\varepsilon)\lambda}{h}\right\rceil.
\]
Both indices tend to infinity and rounding changes either by at most
one. Since $nr_n\to\lambda$,
\[
\frac{r_{n_-}}h\to\frac1{1-\varepsilon}>1,
\qquad
\frac{r_{n_+}}h\to\frac1{1+\varepsilon}<1.
\]
For small $h$, the mean is positive at $n_-$ and strictly negative at
$n_+$. Monotonicity from Appendix~\ref{app:monotonicity} therefore gives
$n_-<n_{{\rm RV},w}\le n_+$. Dividing by $h^{-1}$ bounds the lower and
upper limits of $hn_{{\rm RV},w}$ by $(1-\varepsilon)\lambda$ and
$(1+\varepsilon)\lambda$. Letting $\varepsilon\downarrow0$ yields
$hn_{{\rm RV},w}\to\lambda$. Since $hn_{\rm pop}\to1$, the desired
ratio follows. All dependence on $\theta$ is through $h$, so the result
holds as $|\theta|\to0$ from either sign.

\section{Power sums and reparameterization}

\subsection{Integral comparisons and indexing shifts}\label{app:power-sums}
For fixed $\beta>-1$, integral comparison gives
\[
\sum_{i=2}^{n+1}i^\beta\sim\frac{n^{\beta+1}}{\beta+1}.
\]
For $\beta\ge0$, monotonicity sandwiches the sum between integrals
with endpoints shifted by at most one; the endpoint error is
$O(n^\beta+1)=o(n^{\beta+1})$. For $-1<\beta<0$, the decreasing
summand has a bounded sum--integral discrepancy, also negligible
relative to the divergent integral. Replacing $n+1$ by $n$ in the
leading power is harmless since $(n+1)/n\to1$.

Apply this result with $\beta=\alpha$ and $2\alpha$ to obtain the
total-weight and square-weight sums in Proposition~\ref{prop:power}.
For the reciprocal-weight sum,
\[
\frac{i^\alpha/(i-1)}{i^{\alpha-1}}=\frac{i}{i-1}\to1.
\]
The comparison sum $\sum_{i=2}^{n+1}i^{\alpha-1}$ diverges for
$\alpha>0$. For any fixed tolerance, the relative error in the tail
is eventually within that tolerance, and the finite initial part is
negligible. Thus the reciprocal-weight sum is asymptotic to
$n^\alpha/\alpha$. Dividing by total weight gives
$nr_n\to(\alpha+1)/\alpha$; Theorem~\ref{thm:crossing} then supplies
the crossing ratio.

\subsection{Delay and averaging constants}\label{app:parameterization}
Squaring the total-weight asymptotic and dividing the square-weight
sum by it gives
\[
n\frac{\sum_{i=2}^{n+1}i^{2\alpha}}
{(\sum_{i=2}^{n+1}i^\alpha)^2}
\to\frac{(\alpha+1)^2}{2\alpha+1}.
\]
With $\alpha=1/(\lambda-1)$,
\[
\alpha+1=\frac{\lambda}{\lambda-1},\qquad
2\alpha+1=\frac{\lambda+1}{\lambda-1},
\]
so this constant is $\lambda^2/(\lambda^2-1)$. The oracle result
from Analysis 06 transforms as
\[
\frac{2\alpha^2}{1+\alpha}
=\frac{2}{\lambda(\lambda-1)}.
\]
Both expressions decrease strictly on $\lambda>1$. Their derivatives
are $-2\lambda/(\lambda^2-1)^2$ and
$-2(2\lambda-1)/\{\lambda^2(\lambda-1)^2\}$, respectively.
Their endpoint limits are the ones stated in Section 3.
These are limits of constants obtained for each fixed $\alpha>0$,
not a claim that the asymptotics hold uniformly as $\alpha\downarrow0$.
At $\alpha=0$, Analysis 06 has a vanishing oracle ratio, while
Analysis 09's ordinary-RV crossing has a diverging logarithmic delay
rather than a finite $\lambda$.

\section{Concentration on the moving-crossing scale}

\subsection{Transfer of the Analysis 06 variance}\label{app:transfer}
Write $S_{i-1}=\sum_{t=1}^{i-1}\xi_t$ in this appendix. Expanding
one loss difference gives
\[
(Y_i-\bar Y_{i-1})^2-Y_i^2
=\frac{S_{i-1}^2}{(i-1)^2}
-2\xi_i\frac{S_{i-1}}{i-1}-2\theta\xi_i-\theta^2.
\]
The first two terms are exactly the oracle loss difference used in
Analysis 06. Its score with validation endpoint $n+1$ is therefore
related to the present power-weight score by
\[
\Xi_{n,w}=\Xi_{n+1}^{(\alpha)}
-2\theta\sum_{i=2}^{n+1}i^\alpha\xi_i
-\theta^2\sum_{i=2}^{n+1}i^\alpha.
\]
The endpoint convention differs by one observation, and
$(n+1)^{2\alpha}\sim n^{2\alpha}$ for fixed $\alpha$.

After centering, the oracle score is a quadratic polynomial in the
independent standard Gaussians. Its covariance with each $\xi_i$ is
zero because all centered Gaussian third moments vanish. The constant
bias term has no variance. Consequently the exact variance transfer is
\[
\operatorname{Var}(\Xi_{n,w})
=\operatorname{Var}(\Xi_{n+1}^{(\alpha)})
+4\theta^2\sum_{i=2}^{n+1}i^{2\alpha}.
\]
This is the orthogonality argument of Analysis 09, Appendices D.1--D.2.
Analysis 06's equation (4) and the proof of its proposition, equation
(5), already establish
$\operatorname{Var}(\Xi_{n+1}^{(\alpha)})\sim2n^{2\alpha}/(1+\alpha)$.
We use that result rather than repeating its coefficient calculation.

When $n\theta^2\to c\ge0$, Appendix~\ref{app:power-sums} gives
\[
\frac{\mathbb E\Xi_{n,w}}{n^\alpha}
\to\frac1\alpha-\frac{c}{\alpha+1},\qquad
\frac{\operatorname{Var}(\Xi_{n,w})}{n^{2\alpha}}
\to\frac2{\alpha+1}+\frac{4c}{2\alpha+1}.
\]
For $c\ne\lambda$, the limiting mean coefficient is nonzero, so
the first limit can be expressed with $\sim$ as in Section 4.
At $c=\lambda$, only $\mathbb E\Xi_{n,w}=o(n^\alpha)$ follows;
we do not use an asymptotic equivalence to zero.

\subsection{General lower bound and the boundary}\label{app:general-lower}
The same Gaussian orthogonality works for any deterministic weights,
so \eqref{eq:lower} holds whenever the mean is nonzero. Cauchy--Schwarz
on the $n$ weights gives
\[
(\sum_{i=2}^{n+1}w_i)^2\le n\sum_{i=2}^{n+1}w_i^2,
\qquad
n\frac{\sum_{i=2}^{n+1}w_i^2}{(\sum_{i=2}^{n+1}w_i)^2}\ge1.
\]
For $c\ne\lambda$, $n(r_n-\theta^2)\to\lambda-c\ne0$, so the
mean is eventually nonzero. Rewriting \eqref{eq:lower} yields
\[
\frac{\operatorname{Var}(\Xi_{n,w})}{\{\mathbb E\Xi_{n,w}\}^2}
\ge
\frac{4n\theta^2}{\{n(r_n-\theta^2)\}^2}
\left[n\frac{\sum_{i=2}^{n+1}w_i^2}
{(\sum_{i=2}^{n+1}w_i)^2}\right].
\]
The prefactor tends to $4c/(\lambda-c)^2$, proving
Proposition~\ref{prop:lower}. For $c=\lambda$, its numerator tends to
$4\lambda>0$ and its denominator tends to zero. The bracket is at
least one, so the ratio diverges along any sequence on which it is
defined. At exact mean-zero points, division by the squared mean is
undefined. This proves the boundary conclusion without requiring
a rate for the smaller-order mean.

\subsection{Full power-weight ratio}\label{app:power-ratio}
For $c\ne\lambda$, division of the limits in
Appendix~\ref{app:transfer} gives
\[
\frac{\operatorname{Var}(\Xi_{n,w})}{\{\mathbb E\Xi_{n,w}\}^2}
\to
\frac{2/(\alpha+1)+4c/(2\alpha+1)}
{\{1/\alpha-c/(\alpha+1)\}^2}.
\]
Since $1/\alpha=(\lambda)/(\alpha+1)$, multiplying numerator and
denominator by $(\alpha+1)^2$ makes the denominator
$(\lambda-c)^2$ and the numerator
\[
2(\alpha+1)+\frac{4c(\alpha+1)^2}{2\alpha+1}
=\frac{2\lambda}{\lambda-1}
+\frac{4c\lambda^2}{\lambda^2-1}.
\]
This proves \eqref{eq:full-ratio}, with a finite strictly positive
limit. The boundary follows from Appendix~\ref{app:general-lower};
equivalently, at $c=\lambda$ the variance has a positive
$n^{2\alpha}$ coefficient whereas the squared mean is $o(n^{2\alpha})$.
For $\alpha=1$ and $\lambda=2$, the mean is
$\sim n(1-c/2)$ away from $c=2$, the variance is
$\sim n^2(1+4c/3)$, and the ratio limit is
$(4+16c/3)/(2-c)^2$, agreeing with \eqref{eq:full-ratio}.

\subsection{Comparison of asymptotic regimes}\label{app:regimes}
If $n\theta^2\to0$ with fixed $\alpha>0$, the bias contribution to
the mean and the linear variance are negligible on the scales
$n^\alpha$ and $n^{2\alpha}$. Appendix~\ref{app:transfer} therefore
recovers the Analysis 06 ratio $2\alpha^2/(1+\alpha)$.
For fixed $\lambda>1$, the constant in \eqref{eq:full-ratio} tends
to this value as $c\downarrow0$, but diverges as $c\to\lambda$
because its numerator remains positive while the denominator vanishes.

At fixed $c>0$, $c\ne1$, both numerator terms in
\eqref{eq:full-ratio} diverge as $\lambda\downarrow1$, while
$(\lambda-c)^2\to(1-c)^2>0$. This recovers the qualitative
obstruction at Analysis 09's alignment boundary. At $c=1$, the
denominator also vanishes and the constant again diverges.
As $\lambda\to\infty$ with $c>0$ fixed, the numerator tends to
$2+4c$ and the denominator grows as $\lambda^2$, so the constant
tends to zero. In this comparison $n/n_{{\rm RV},w}\to c/\lambda$
in the preceding fixed-parameter limit, placing the evaluation
increasingly far before the weighted crossing.

All parameter comparisons follow the fixed-$\alpha$, fixed-$c$
asymptotics; no uniform or simultaneous limit is asserted. In
particular, neither the small oracle benchmark nor a bounded
square-weight factor gives a vanishing ratio for a fixed finite
$\lambda$ and fixed $c>0$. Nor does a nonvanishing Chebyshev ratio
by itself establish a positive limiting wrong-selection probability.

\end{document}
